\scanpage{215}
\mspage{70}

au même de dire que $R(p)$ est un \uline{ouvert} de $Q$ stable par Sup quelconques (ce qui implique qu'il contient $\varnothing_Q$). Dualement, pour tout $q\in Q$, $\,{}^{\unsure{s|t}}\!R(q)=\{p\in P\mid(p,q)\in R\}$ est une partie \uline{ouverte} de $P$, stable par Sup quelconques, i.e. une partie \unsure{ouverte|ouvert} \deleted{filtrante} ayant un plus grand élément.

En d'autres termes, quand on postule la condition (D) pour $A$ resp. $B$ réduit à un élément, alors \unsure{elle} reste vérifiée pour \unsure{tte|toute} partie. Cela provient des \deleted{\unsure{de|la}} formules

\[
\left\{\begin{aligned}
 \left(\bigcup_i A_i\right)^{\unsure{\cup|\smile}}&=\bigcap_i A_i^\cup\\
 \left(\bigcup_j B_j\right)^{\unsure{\cap|\frown}}&=\bigcap_j B_j^\cap,
\end{aligned}\right.
\tag*{(5.2.)}
\]

et du fait que l'intersection de parties ouvertes de la forme $P_{/\alpha_i}$ de $P$ est de la forme $P_{/\alpha}$, où $\alpha=\mathop{\mathrm{Inf}}_i\alpha_i$, et de $=$ dans $Q$. Ainsi on trouve

\[
\left\{\begin{aligned}
 \beta\left(\bigcup_i A_i\right)&=\mathop{\mathrm{Inf}}_i\beta(A_i)\\
 \alpha\left(\bigcup_j B_j\right)&=\mathop{\mathrm{Inf}}_j\alpha(B_j).
\end{aligned}\right.
\tag*{(5.3.)}
\]

On posera, pour une partie $A$ de $P$ resp. $B$ de $Q$

\[
 A^{\unsure{\sharp}}=(A^\cup)^\cap,\qquad B^\flat=(B^\cap)^\cup
\tag*{(5.4).}
\]

\deleted{\unsure{(5.5)}} de sorte qu'on aura

\[
 A\subset A^\sharp,\qquad B\subset B^\flat
\tag*{(5.5.)}
\]

\deleted{On aura, pour tt $A$ resp.}

\deleted{\[
 (A^\cup)
 \tag*{(5.6)}
\]}

On dira que $A$ est \uline{$R$-clos} si $A=A^\sharp$ \inserted{(ou encore si $A$ est de la forme $B'^\cap$ pour $B'\subset Q$, cf 5.6)}, que
% archive sheet number: 215
